\documentclass[11pt]{article}
\usepackage[a4paper,margin=25mm]{geometry}
\usepackage{amsmath,amssymb,amsthm,mathtools}
\usepackage[T1]{fontenc}
\usepackage{lmodern}
\usepackage[colorlinks=true,linkcolor=blue]{hyperref}
\newtheorem{theorem}{Theorem}
\newtheorem{lemma}[theorem]{Lemma}
\newtheorem{proposition}[theorem]{Proposition}
\newcommand{\E}{\mathcal E}
\newcommand{\PP}{\mathcal P}
\title{The four-colour component optimum}
\author{}
\date{}
\begin{document}
\maketitle

We determine the optimum of the finite component program below by an
elementary analytic argument. The main step is to exclude a large atom before
estimating the remaining four-atom problem by its moments.
All logarithms are natural.

\section{The program and its optimum}

Choose a finite set of atoms $[m]$, masses $a_i>0$ with $\sum_i a_i=1$,
and four partitions $\PP_c$ of $[m]$. Choose symmetric densities
$0\le p^c_{ij}<1$ such that
\[
 \sum_{c=1}^4p^c_{ij}=1,
 \qquad p^c_{ij}=0\quad\text{if $i,j$ lie in different parts of $\PP_c$}.
\]
For $C\in\PP_c$, define
\[
 \E_{c,C}=-\sum_{i,j\in C}a_i a_j\log(1-p^c_{ij}),
 \qquad t=\max_{c,C}\E_{c,C}.
\]
The sums include ordered pairs and diagonal terms. Let $\beta_4$ be the
infimum of $\sqrt t$ over this program.

\begin{theorem}\label{thm:main}
Let $p\in(0,1)$ be the unique solution of $3p^6+2p-2=0$. Then
\[
 \beta_4^2=t_*=-\frac34\log p.
\]
The optimum is attained by four equal masses, the partitions
\[
 \{12,34\},\qquad\{13,24\},\qquad\{14,23\},\qquad\{1234\},
\]
diagonal densities $(1/3,1/3,1/3,0)$, and, on each distinct pair,
density $p$ in its pairing colour and $1-p$ in the fourth colour.
\end{theorem}

Here $\{12,34\}$ means $\{\{1,2\},\{3,4\}\}$, and similarly for the
other partitions. The fourth component in the construction has energy
$-\frac34\log p$, while each of the six paired components has energy
\[
 \frac18\log\frac{3}{2(1-p)}=-\frac34\log p.
\]
Thus the construction gives the upper bound. We prove the lower bound below.
Throughout, put
\begin{equation}\label{eq:parameters}
 w=\frac{1-p}{p},\qquad
 w(1+w)^5=\frac32,\qquad
 T=\frac{2213}{10000}.
\end{equation}
We will use $t_*<T$. Here is an explicit elementary verification. For
$p_0=1489/2000$, direct rational comparison gives
$3p_0^6+2p_0-2<0$, so $p>p_0$. With $y=511/3489=(1-p_0)/(1+p_0)$,
the positive-term logarithm series gives
\begin{equation}\label{eq:threshold}
 t_*<\frac34\log\frac1{p_0}
 <\frac32\left(y+\frac{y^3}{3}+\frac{y^5}{5(1-y^2)}\right)
 <\frac{2213}{10000}.
\end{equation}
The last comparison is rational.

\section{The possible component patterns}

Label each atom by its four component labels. Any two atoms agree in at least
two coordinates: otherwise at most one colour is available between them,
forcing a density equal to $1$. Atoms with identical labels can be merged,
using mass-weighted average densities, including ordered pairs on a merged
diagonal. Convexity of $-\log(1-x)$ shows that no component energy increases.
It therefore suffices to consider distinct words of length four whose pairwise
Hamming distances are at most two.

\begin{lemma}\label{lem:patterns}
After relabelling coordinates and their symbols, every such family is
contained in one of the following:
\begin{enumerate}
\item a family with two coordinates fixed;
\item a radius-one Hamming ball;
\item the four-word family $0000,0110,1010,1100$.
\end{enumerate}
The centre of the ball need not belong to the family.
\end{lemma}
\begin{proof}
Choose a word and relabel it $0000$. All other supports have size at most
two. The two-element supports are pairwise intersecting, so they either have
a common coordinate or form a triangle.

If there are no two-element supports, the family lies in the ball about
$0000$. If there is exactly one distinct two-element support, every singleton
support is contained in it, giving two fixed coordinates.

Suppose there are at least two distinct two-element supports and they have
a common coordinate. Their nonzero symbols in that coordinate must coincide,
since otherwise their distance would be three. Comparison with a different
support gives the same conclusion for every word on every such support.
No singleton support outside the common coordinate is compatible with all
these words. Changing that coordinate of the centre to the common nonzero
symbol puts the entire family in a radius-one ball.

Finally, if the supports form a triangle, comparison of two intersecting
supports fixes one nonzero symbol in each of the three coordinates.
There is one possible word on each support, and no singleton support is
compatible with all three. This gives the stated four-word family.
\end{proof}

The third pattern is the three-pairing pattern in Theorem~\ref{thm:main}.
We first treat it, allowing zero masses to include its subfamilies.

\section{The four-atom pattern}

Write $E_0$ for the energy of the spanning colour, and $E_{ij}$ for that
of the paired component on $\{i,j\}$. Assume $a_i\ge0$ and
$\sum_{i=1}^4a_i=1$. Define
\[
 H(x)=(1+x)\log(1+x)-x\log x\qquad(x>0).
\]
For two available colours, minimizing their weighted logarithmic cost gives
\begin{equation}\label{eq:two-colour}
 -\log(1-q_0)-x\log(1-q_1)\ge H(x)
 \qquad(q_0+q_1=1).
\end{equation}

\begin{lemma}\label{lem:large-atom}
If all seven component energies are at most $t_*$, then
$\max_i a_i<2/5$.
\end{lemma}
\begin{proof}
Fix an atom of mass $x$. Give the spanning energy weight $1$ and the three
paired energies incident to this atom weight $1/2$ each. On its diagonal,
the tangent inequality at densities $(0,1/3,1/3,1/3)$ gives weighted
logarithmic cost at least $(3/2)\log(3/2)$: the tangent slopes are $1$ in
the spanning colour and $3/4$ in each of the other three colours.
For every pair from this atom to another atom, \eqref{eq:two-colour}
gives cost at least $H(1/2)$. Dropping all other nonnegative contributions,
\[
 \frac52t_*\ge f(x):=
 \frac32\log\frac32\,x^2+2H(1/2)x(1-x).
\]
This quadratic is concave. The elementary estimates
$\log(3/2)>81/200$ and $\log2>69/100$ give
\[
 f(2/5)=\frac{24}{25}\log\frac32+\frac6{25}\log2
 >\frac{693}{1250}>\frac52T,
 \qquad
 f(1)>\frac{243}{400}>\frac52T.
\]
Both logarithm estimates follow by retaining two terms in
\begin{equation}\label{eq:logseries}
 \log\frac{1+z}{1-z}
 =2\sum_{j=0}^{\infty}\frac{z^{2j+1}}{2j+1}\qquad(0<z<1).
\end{equation}
Concavity and $t_*<T$ now exclude every $x\in[2/5,1]$.
\end{proof}

We next record the elementary bounds needed for the four-atom inequality.
Put
\[
 v=3w-\frac14,\qquad b=\log\frac32,
 \qquad h=H(w),\qquad h'=H'(w)=\log\frac{1+w}{w}.
\]

\begin{lemma}\label{lem:entropy}
The following inequalities hold:
\begin{gather}
 \frac13<w<\frac{11}{32},\qquad v>\frac34,\label{eq:wbound}\\
 \alpha:=vb-h+\frac54h'>\frac54,
 \qquad \gamma:=b-h'>-1,\label{eq:coefficients}\\
 H(w+d)\ge h+h'd-2d^2\qquad\left(|d|\le\frac3{20}\right).
 \label{eq:Hbound}
\end{gather}
\end{lemma}
\begin{proof}
The function $x(1+x)^5$ increases. At $x=1/3$ its value is
$1024/729<3/2$, while at $x=11/32$ its value exceeds $3/2$ by
$6480137/1073741824$. This proves \eqref{eq:wbound}.
Also
\[
 \gamma=\log\frac{3w}{2(1+w)}>\log\frac38>-1,
\]
since $e>1+1+1/2+1/6=8/3$.

For $A(x)=(3x-1/4)b-H(x)+(5/4)H'(x)$, we have
\[
 A'(x)=3b-H'(x)-\frac{5}{4x(1+x)}<0
 \quad\left(\frac13\le x\le\frac{11}{32}\right),
\]
using $b<1/2$, $H'(x)>0$, and $x(1+x)<1/2$. Moreover,
\[
 32A(11/32)=25\log\frac32+29\log\frac{32}{11}
                  -3\log\frac{43}{32}.
\]
The series \eqref{eq:logseries} and $\log(1+x)<x$ give
\[
 \log\frac32>\frac{152}{375},\qquad
 \log\frac{32}{11}=\log3-\log\frac{33}{32}
 >1+\frac1{12}+\frac1{80}-\frac1{32}=\frac{511}{480}.
\]
Also $\log(43/32)<3/10$, since
$e^{3/10}>1+3/10+9/200=269/200>43/32$. Hence
\[
 32\alpha>25\frac{152}{375}+29\frac{511}{480}-\frac9{10}
 =\frac{6417}{160}>40.
\]

Finally, $H''(x)=-1/[x(1+x)]$. For $d\ne0$, Taylor's integral formula
gives
\[
 \frac{h+h'd-H(w+d)}{d^2}
 =\int_0^1\frac{1-s}{(w+sd)(1+w+sd)}\,ds.
\]
Here $w>1/3$ and $w+d>11/60$. Convexity of $1/x$ therefore bounds
the integrand by
$(1-s)\bigl(3(1-s)+(60/11)s\bigr)$.
Its integral is $21/11<2$, proving \eqref{eq:Hbound}.
\end{proof}

Set
\[
 w_{ij}=w+\frac{a_i+a_j-1/2}{2},\qquad
 S_i=\sum_{j\ne i}w_{ij}=v+a_i,
 \qquad
 D_i=\sum_{j\ne i}w_{ij}\log\frac{S_i}{2w_{ij}}.
\]
These weights are positive, $\sum_{i<j}w_{ij}=6w$, and $S_i<2$.
By \eqref{eq:two-colour}, each off-diagonal pair contributes at least
$H(w_{ij})$ to the weighted logarithmic cost. On diagonal $i$, use
formal densities
\[
 r_0=0,\qquad r_{ij}=1-\frac{2w_{ij}}{S_i}.
\]
They sum to $1$ and are all less than $1$, although some may be negative.
The function $-\log(1-x)$ is convex on the whole interval $x<1$.
Its weighted tangent slopes at these densities are $1$ in the spanning
colour and $S_i/2$ in the three paired colours. At feasible densities
$q_0,q_{ij}\ge0$, the tangent bound is
$D_i+(1-S_i/2)q_0\ge D_i$. Consequently
\begin{equation}\label{eq:dual}
 E_0+\sum_{i<j}w_{ij}E_{ij}\ge N(a),
 \qquad
 N(a):=\sum_i a_i^2D_i+2\sum_{i<j}a_i a_jH(w_{ij}).
\end{equation}

\begin{proposition}\label{prop:four-atoms}
If $0\le a_i\le2/5$ and $\sum_i a_i=1$, then
\begin{equation}\label{eq:stability}
 N(a)\ge(1+6w)t_*+\frac3{100}\sum_i(a_i-1/4)^2.
\end{equation}
In particular, every feasible four-atom pairing pattern has $t\ge t_*$.
\end{proposition}
\begin{proof}
Write
\[
 q=\sum_i a_i^2,\qquad r=\sum_i a_i^3,\qquad s=\sum_i a_i^4.
\]
Using $\log x\le x-1$, we have
\begin{align*}
 D_i
 &=S_i b-\sum_{j\ne i}w_{ij}\log\frac{3w_{ij}}{S_i}\\
 &\ge S_i b-\frac3{S_i}\sum_{j\ne i}(w_{ij}-S_i/3)^2\\
 &=S_i b-\frac{3q-4a_i^2+2a_i-1}{4S_i}.
\end{align*}
The numerator in the last term is nonnegative, being twelve times the
preceding sum of squares. Since $S_i\ge3/4$, it follows that
\begin{equation}\label{eq:diagonal-moments}
 \sum_i a_i^2D_i\ge b(vq+r)-\frac{3q^2-4s+2r-q}{3}.
\end{equation}

The mass restriction implies $1/5\le a_i+a_j\le4/5$, so
$d_{ij}=(a_i+a_j-1/2)/2$ lies in $[-3/20,3/20]$.
Expanding the sums in \eqref{eq:Hbound} gives
\begin{align}
 2\sum_{i<j}a_i a_jH(w_{ij})
 &\ge h(1-q)+h'\left(\frac54q-r-\frac14\right)\notag\\
 &\quad-\left(2r+q^2-2s-\frac98q+\frac18\right).
 \label{eq:offdiag-moments}
\end{align}
Indeed, the quadratic loss is
$\sum_{i<j}a_i a_j(a_i+a_j-1/2)^2$, which is the last parenthesis.

At $a_i=1/4$ we have
\begin{equation}\label{eq:uniform}
 N_0:=N(1/4,1/4,1/4,1/4)=\frac34(wb+h)=(1+6w)t_*;
\end{equation}
the last identity uses $\log(3/2)=\log w+5\log(1+w)$.
Subtracting this equality from \eqref{eq:diagonal-moments} and
\eqref{eq:offdiag-moments}, and using $q\ge1/4$, $r\ge1/16$ and
\eqref{eq:coefficients}, yields
\begin{align*}
 N(a)-N_0
 &\ge \alpha(q-1/4)+\gamma(r-1/16)
       -\left(2r+q^2-2s-\frac98q+\frac18\right)
       -\frac{3q^2-4s+2r-q}{3}\\
 &\ge \frac{65}{24}q-\frac{11}{3}r+\frac{10}{3}s-2q^2-\frac38.
\end{align*}
To estimate the last polynomial, put
\[
 x_i=a_i-\frac14,\qquad
 Q=\sum_i x_i^2,\qquad R=\sum_i x_i^3,\qquad U=\sum_i x_i^4.
\]
Since $\sum_i x_i=0$, the polynomial is exactly
\[
 \frac5{24}Q-\frac13R+\frac{10}{3}U-2Q^2.
\]
Now $x_i\le3/20$ gives $R\le(3/20)Q$, and Cauchy--Schwarz gives
$U\ge Q^2/4$. Also
\[
 q\le\left(\frac25\right)^2+\left(\frac25\right)^2
            +\left(\frac15\right)^2=\frac9{25},
 \qquad Q\le\frac{11}{100}.
\]
For this bound on $q$, move mass from a smaller positive coordinate to a
larger coordinate below $2/5$ until one reaches $0$ or $2/5$; this never
decreases the sum of squares and ends at a permutation of
$(2/5,2/5,1/5,0)$. Therefore
\[
 N(a)-N_0\ge\frac{19}{120}Q-\frac76Q^2
 \ge\left(\frac{19}{120}-\frac76\frac{11}{100}\right)Q
 =\frac3{100}Q.
\]
This proves \eqref{eq:stability}. If a pairing pattern had $t<t_*$,
Lemma~\ref{lem:large-atom} would put it in this mass range, while
\eqref{eq:dual} would give
$(1+6w)t\ge N(a)\ge(1+6w)t_*$, a contradiction.
\end{proof}

\section{Bounds for the other component patterns}

For $t>0$ and $s\ge0$, define
\[
 F_t(s)=s^2(1-e^{-t/s^2}),\quad F_t(0)=0,
 \qquad
 g_t(s)=3s^2\log\frac{3-e^{-t/s^2}}2,\quad g_t(0)=0.
\]
A component of mass $s$ and energy at most $t$ has total ordered edge mass
at most $F_t(s)$, by Jensen's inequality. The function $F_t(s)$ increases
with $s$; differentiating reduces this to $e^z\ge1+z$.

\begin{lemma}\label{lem:capacity}
For every $t>0$ and $s\ge0$,
\[
 F_t(s)\le\frac{16}{25}\sqrt t\,s,
 \qquad
 g_t(s)\le\frac{413}{500}\sqrt t\,s.
\]
\end{lemma}
\begin{proof}
For the first inequality, maximize $(1-e^{-z})/\sqrt z$.
Its unique positive critical point $z_0$ satisfies $e^{z_0}=1+2z_0$.
The elementary Taylor bounds
\[
 e<\frac83+\frac{1/24}{1-1/5}=\frac{87}{32},
 \qquad
 e^{1/4}<\frac54+\frac{1/32}{1-1/12}=\frac{113}{88}
\]
give $e^{5/4}<9831/2816<7/2$, so $z_0>5/4$.
Since $2\sqrt z/(1+2z)$ decreases for $z>1/2$, the maximum is
\[
 \frac{2\sqrt{z_0}}{1+2z_0}<\frac{2\sqrt5}{7}<\frac{16}{25}.
\]
Substitute $z=t/s^2$.

For the second inequality, maximize
$3\log((3-e^{-z})/2)/\sqrt z$. The derivative has the sign of
\[
 G(z)=\frac{2z}{3e^z-1}-\log\frac{3-e^{-z}}2,
 \qquad
 G'(z)=\frac{3e^z(1-2z)-1}{(3e^z-1)^2}.
\]
The numerator of $G'$ decreases strictly from $2$. Thus $G$ first
increases and then decreases; since $G(0)=0$ and
$G(z)\to-\log(3/2)$, it has one positive zero $z_1$.

To locate it sufficiently, put $x=41/40$. Taylor's series gives
\[
 \frac{2787}{1000}
 <\sum_{j=0}^{8}\frac{x^j}{j!}
 <e^x
 <\sum_{j=0}^{8}\frac{x^j}{j!}
       +\frac{x^9}{9!(1-x/10)}
 <\frac{27871}{10000}.
\]
For $y=1391/5000$, direct rational comparison gives
\[
 1+y+\frac{y^2}{2}+\frac{y^3}{6}+\frac{y^4}{24}
 >\frac{13207}{10000}>\frac{73613}{55742}
 =\frac{3-10000/27871}{2}.
\]
Consequently
\[
 G(41/40)>\frac{20500}{73613}-\frac{1391}{5000}>0,
\]
so $z_1>41/40$. At this zero the maximum quotient equals
$6\sqrt{z_1}/(3e^{z_1}-1)$. This expression decreases for $z\ge1/2$.
Using $\sqrt{41/40}<81/80$, we obtain
\[
 \frac{6\sqrt{z_1}}{3e^{z_1}-1}
 <\frac{6(81/80)}{3(2787/1000)-1}
 =\frac{6075}{7361}<\frac{413}{500}.
\]
Again substitute $z=t/s^2$.
\end{proof}

We exclude the first two patterns of Lemma~\ref{lem:patterns} even at
the threshold $T>t_*$. To keep the comparisons explicit, set
\[
 P=T-\frac{T^2}{2}+\frac{T^3}{6}.
\]
Taylor's inequality and squaring give
\begin{equation}\label{eq:Tbounds}
 F_T(1)<P,\qquad \frac{47}{100}<\sqrt T<\frac{941}{2000}.
\end{equation}

If two coordinates are fixed, two colours have spanning components.
Each of the other two colour partitions has total component mass $1$.
If $t\le T$, summing capacities and using Lemma~\ref{lem:capacity} gives
\[
 1\le2F_T(1)+\frac{32}{25}\sqrt T
 <2P+\frac{32}{25}\frac{941}{2000}<1,
\]
a contradiction; the last inequality is rational.

It remains to exclude a radius-one Hamming ball. Let $z$ be the mass of
the centre, possibly zero, and let $x_c$ be the mass on arm $c$.
Put $S=\sum_c x_c=1-z$ and $Q=\sum_c x_c^2$. Colour $c$ has a main
component of mass $1-x_c$ and a singleton component on each atom of
arm $c$. Write $a$ for the individual arm-atom masses, so $\sum_a a=S$.
Set
\[
 A=2\log2,\qquad B=3\log(3/2),\qquad C=4\log(4/3).
\]
For an ordered pair in different arms, exactly two main colours are
available, giving total logarithmic cost at least $A$. For distinct atoms
in the same arm, or a centre--arm pair, three main colours are available,
giving cost at least $B$. The centre diagonal gives cost at least $C$.
On the diagonal of an arm atom of mass $a$, its singleton-colour density
is at most $1-e^{-t/a^2}$. The three main-colour densities therefore sum
to at least $e^{-t/a^2}$, giving main-colour cost at least
$-3\log(1-e^{-t/a^2}/3)$. Thus the sum $E$ of the four main-component
energies satisfies
\begin{equation}\label{eq:ball}
 E\ge AS^2+(B-A)Q+2BzS+Cz^2-\sum_a g_t(a).
\end{equation}

Suppose $t\le T$, and put $K=(413/500)\sqrt T$. By
Lemma~\ref{lem:capacity}, the part of the right side of \eqref{eq:ball}
other than $(B-A)Q$ is at least
\[
 C+(2B-2C-K)S+(A-2B+C)S^2.
\]
This quadratic decreases on $[0,1]$: its quadratic coefficient is
$A-2B+C=\log(65536/59049)>0$, while its derivative at $1$ satisfies
\[
 2(A-B)-K=2\log\frac{32}{27}-K
 <\frac{10}{27}-\frac{413}{500}\frac{47}{100}<0.
\]
Its minimum is $A-K$, and $E\le4T$ therefore implies
\[
 Q\ge\frac{A-K-4T}{A-B}>\frac{33}{50}.
\]
For the strict comparison, four terms of \eqref{eq:logseries} give
$\log2>6931/10000$, and three give $\log(3/2)>4054/10000$. Hence
\[
 \frac{17}{50}A+\frac{33}{50}B>\frac{637}{500}
 >\frac{413}{500}\frac{941}{2000}+4T>K+4T,
\]
which is exactly the required inequality.

Since $Q\le(\max_c x_c)S\le\max_c x_c$, some main component has
mass less than $17/50$. Summing the capacities of all main and singleton
components now gives
\[
 1\le3F_T(1)+F_T(17/50)+\frac{16}{25}\sqrt T.
\]
But $T/(17/50)^2<2$, and the bound $e<87/32$ used above gives
$e^2<15/2$. Thus
\[
 F_T(17/50)<\frac{13}{15}\left(\frac{17}{50}\right)^2,
\]
and the preceding capacity sum is strictly less than
\[
 3P+\frac{13}{15}\left(\frac{17}{50}\right)^2
       +\frac{16}{25}\frac{941}{2000}<1.
\]
This final comparison is rational, giving a contradiction.

The first two component patterns cannot occur with $t\le t_*$, and
Proposition~\ref{prop:four-atoms} gives $t\ge t_*$ in the third.
Together with the attaining construction, this proves
Theorem~\ref{thm:main}.

\end{document}
